\section{Introduction}

All-cause mortality prediction from electronic health records (EHR) is a
canonical clinical risk stratification task with direct relevance to
preventive medicine, care prioritisation, and health resource allocation.
Unlike typical classification problems, mortality prediction requires models
to reason over \emph{structured multimodal data}: longitudinal sequences of
disease events with precise temporal ordering, numeric biomarker measurements
with clinical significance thresholds, and demographic context.
The underlying structure --- temporal dependencies, numeric precision, causal
relationships between comorbidities --- is what makes this task hard and what
makes it a meaningful testbed for structured multimodal intelligence.

Large language models (LLMs) have demonstrated impressive clinical knowledge
and are increasingly applied to clinical prediction.
Two distinct paradigms have emerged: \emph{zero-shot autoregressive inference},
where the LLM directly generates survival-related outputs from tokenised EHR
sequences~\cite{delphi2024}, and \emph{supervised embedding}, where a frozen
LLM encodes patient records as dense vectors combined with a lightweight
survival head.
Both paradigms make implicit claims about structural grounding, yet the
relative importance of specific structural properties --- temporal ordering,
numeric precision, encoder choice, and fusion method --- has not been
systematically evaluated under controlled conditions.

We address this gap with a three-axis controlled study on UK Biobank
($n\,=\,124{,}158$), the largest such benchmark for LLM-based mortality
prediction.
We vary:
\textbf{(A)} \emph{input representation} across five temporal-encoding
settings, from cross-sectional bag-of-features to fully unified clinical
prose;
\textbf{(B)} \emph{encoder architecture} (Qwen3-8B vs.\
Bio\_ClinicalBERT) for the three embedding-based settings, to test whether
findings are encoder-specific;
\textbf{(C)} \emph{feature fusion method} (concatenation vs.\ element-wise
summation) for settings that combine embeddings with structured features.
All embedding variants use a uniform CoxPH survival head, enabling
apples-to-apples comparison of the structural encoding choices.

Our main contributions are:
\begin{itemize}
  \item A systematic, reproducible three-axis benchmark disentangling
        representation, encoder, and fusion effects for LLM-based mortality
        prediction on UK Biobank.

  \item Evidence that \emph{temporal ordering} is the key structural property:
        decoupled age--event encoding (Setting~4) achieves a $+$0.008 C-index
        advantage over prose serialisation of identical events, holding encoder,
        raw features, and fusion fixed; bootstrap CIs will determine significance.

  \item The finding that Qwen3-8B can encode numeric biomarker structure
        from prose: unified summary embedding (Setting 5) with \emph{no raw
        features} matches a 39-feature CoxPH baseline, demonstrating
        implicit numeric precision in LLM embeddings.

  \item Evidence that zero-shot Delphi (Setting 2) has a structural
        calibration gap: poor near-term AUC and high Brier score reveal
        that autoregressive inference without task-specific structural
        grounding is insufficient for calibrated risk stratification.

  \item A \emph{planned} encoder-generalisation comparison
        (Qwen3-Embedding (8B) vs.\ Bio\_ClinicalBERT) and feature-fusion
        comparison (concatenation vs.\ summation) that will characterise how much
        the structural encoding choice --- rather than model family or fusion
        method --- drives prediction quality.
\end{itemize}
